\documentclass{article}

\usepackage[english]{babel}

\usepackage[a4paper,top=2cm,bottom=2cm,left=3cm,right=3cm,marginparwidth=1.75cm]{geometry}
\usepackage{amsmath}
\usepackage{unicode-math}
\usepackage{graphicx}
\usepackage[colorlinks=true, allcolors=blue]{hyperref}
\usepackage{listings}
\usepackage{minted}
\usepackage{booktabs}
\usepackage{physics}
\usepackage{amssymb}%braket schreibweise
\setlength{\parindent}{0pt}

\title{Physics of Complex Systems: Exercise 10}
\author{Leopold Schaller and Daniel Lender}

\begin{document}
\maketitle

\section*{11.1: Entanglement of a pure 2-qubit state}
\subsection*{(a)}

We calculated the matrix and verified it's purity (the trace of it's square has to be one) in python with this code:

\begin{minted}{python}
import sympy as sp
i = sp.I
sqrt2 = sp.sqrt(2)

psi = sp.Matrix([
    sp.Rational(1, 5),
    (2 * i * sqrt2) / 5,
    sp.Rational(-4, 15),
    (-8 * i * sqrt2) / 15])

rho = psi * psi.H

print("latex code")
print(sp.latex(rho))

rho_squared = rho**2
purity = rho_squared.trace()

print("purity")
print(f"Tr(rho²) = {purity}")
\end{minted}

\begin{align}
\rho=
    \begin{pmatrix}
    \frac{1}{25} & - \frac{2 \sqrt{2} i}{25} & - \frac{4}{75} & \frac{8 \sqrt{2} i}{75}\\[0.1cm]
    \frac{2 \sqrt{2} i}{25} & \frac{8}{25} & - \frac{8 \sqrt{2} i}{75} & - \frac{32}{75}\\[0.1cm]
    - \frac{4}{75} & \frac{8 \sqrt{2} i}{75} & \frac{16}{225} & - \frac{32 \sqrt{2} i}{225}\\[0.1cm]
    - \frac{8 \sqrt{2} i}{75} & - \frac{32}{75} & \frac{32 \sqrt{2} i}{225} & \frac{128}{225}
    \end{pmatrix}
\end{align}

\begin{align}
    \text{Tr}(\rho^2)=1
\end{align}
\subsection*{(b)}

We calculated the reduced matrices with the by carrying out the partial traces in python:
\begin{minted}{python}
import sympy as sp

i = sp.I
sqrt2 = sp.sqrt(2)

psi = sp.Matrix([
    sp.Rational(1, 5),
    (2 * i * sqrt2) / 5,
    sp.Rational(-4, 15),
    (-8 * i * sqrt2) / 15
])
rho = psi * psi.H

rho_L = sp.Matrix([
    [rho[0,0] + rho[1,1], rho[0,2] + rho[1,3]],
    [rho[2,0] + rho[3,1], rho[2,2] + rho[3,3]]
])

rho_R = sp.Matrix([
    [rho[0,0] + rho[2,2], rho[0,1] + rho[2,3]],
    [rho[1,0] + rho[3,2], rho[1,1] + rho[3,3]]
])

print("rho_L")
print(sp.latex(rho_L))
print("\n rho_R")
print(sp.latex(rho_R))
\end{minted}

\begin{align}
    \rho_L=
    \begin{pmatrix}
    \frac{9}{25} & - \frac{12}{25}\\
    - \frac{12}{25} & \frac{16}{25}
    \end{pmatrix} \hspace{1cm}
    \rho_R=
    \begin{pmatrix}\frac{1}{9} & - \frac{2 \sqrt{2} i}{9}\\\frac{2 \sqrt{2} i}{9} & \frac{8}{9}\end{pmatrix}
\end{align}

\subsection*{(c)}
A pure state is called entangled if the reduced density matrices are mixed. To check if $\rho$ is entangle we simply check, if the reduced density matrices are pure.
With Python (code similar to the code above) we calculated:
\begin{align}
        \text{Tr}(\rho_L^2)=1
        \text{Tr}(\rho_R^2)=1
\end{align}
This means, $\rho$ is not entangled.

\section*{11.2: Quantum Cryptography}
\subsection*{(a)}

\subsection*{(b)}

\subsection*{(c)}

\subsection*{(d)}

\subsection*{(e)}


\end{document}
# lualatex -shell-escape exercise11_schaller_lender.tex
